% !TEX program = xelatex
% =============================================================================
% 模板 / TEMPLATE —— 使用可复用的 sustech 主题
%
% 用法：复制本文件为 main.tex，替换所有 〈…〉 占位符为真实内容即可。
% 构建：latexmk -xelatex main_template.tex
%       （latexmkrc 已把 ./sustech-theme// 加入 TEXINPUTS）
%
% 本模板演示主题提供的全部能力：
%   配色 / 章节分隔页(含一句话主旨) / \shl \keyword \brandemph \hlbox
%   \metric / callout / block / exampleblock / 表格(\theadrow\altrow\thc)
%   \fitfigure + \figcap / 参考文献 / 附录 backup / 标题页 logo
% =============================================================================
\documentclass[aspectratio=169,10pt]{ctexbeamer}

\usetheme{sustech}        % Madrid 基础 + 学术配色 + 排版宏（见 sustech-theme/）

\usepackage{amssymb}
\usepackage{tikz}
\graphicspath{{figures/}}

% --- 占位辅助（仅模板用；真实稿件可删除）------------------------------------
% 行内占位文字
\newcommand{\ph}[1]{\textcolor{sustechgrey}{〈#1〉}}
% 图片占位框（替换为 \fitfigure{your_image} 即可）
\newcommand{\phfig}{%
  \fcolorbox{cardgrey}{bgsoft}{%
    \parbox[c][0.42\textheight][c]{0.96\linewidth}{\centering
      \color{sustechgrey}（图片占位 \textbar\ \texttt{\textbackslash fitfigure\{...\}}）}}}

% --- 元信息（替换为你的题目 / 作者）------------------------------------------
\title[\ph{短标题}]{\ph{中文主标题}}
\subtitle{\ph{英文副标题 / English subtitle}}
\author[\ph{汇报人}]{\ph{汇报人姓名}}
\institute[\ph{机构}]{\ph{机构} \textperiodcentered\ \ph{会议 / 期刊}}
\date{\ph{YYYY-MM-DD} \textperiodcentered\ \ph{地点}}
% 标题页 logo 由 sustech 主题自动提供；如需更换：\renewcommand{\sustechlogo}{...}

\begin{document}

% ============================ 标题页 ============================
\begin{frame}[plain]
  \titlepage
\end{frame}

% ============================ 目录 ============================
\begin{frame}{目录}
  \tableofcontents[hideallsubsections]
\end{frame}

% ============================ 1. 作者介绍 ============================
\renewcommand{\secblurb}{\ph{本节一句话主旨：团队 / 作者与本文的关系}}
\section{作者介绍}

\begin{frame}{作者介绍}
  \begin{columns}[T]
    \begin{column}{\zonehalf}
      \begin{block}{核心作者}
        \begin{itemize}
          \item \shl{\ph{一作姓名}}：\ph{单位、研究方向、与本文的角色}。
          \item \shl{\ph{合作者}}：\ph{贡献}。
        \end{itemize}
      \end{block}
    \end{column}
    \begin{column}{\zonehalf}
      \begin{block}{通讯作者}
        \begin{itemize}
          \item \shl{\ph{通讯作者}}：\keyword{\ph{领域}} \ph{代表性背景}。
          \item \ph{实验室 / 代表性工作}。
        \end{itemize}
      \end{block}
    \end{column}
  \end{columns}
  \tightgap
  \begin{callout}[研究脉络]
    \ph{用一句话概括该团队的研究主线与本文的位置}。
  \end{callout}
\end{frame}

% ============================ 2. 研究背景与动机 ============================
\renewcommand{\secblurb}{\ph{本节一句话主旨：本文要解决的问题与核心贡献}}
\section{研究背景与动机}

\begin{frame}{这篇论文在做什么？}
  \begin{columns}[T]
    \begin{column}{\zonetext}
      \begin{itemize}
        \item 提出 \shl{\ph{方法名}}：一个\keyword{\ph{关键特性}}的\ph{系统 / 方法}。
        \item 把 \brandemph{\ph{要素一}}、\brandemph{\ph{要素二}} 整合在一起。
        \item 覆盖 \ph{任务 / 场景范围}。
        \item 仅需 \hlbox{\ph{关键数字}} 即可 \keyword{\ph{核心收益}}。
      \end{itemize}
    \end{column}
    \begin{column}{\zonefig}
      \centering
      \phfig
      \tightgap
      \figcap{1}{\ph{图注}}
    \end{column}
  \end{columns}
  \deckgap
  \centering
  \metric{\ph{指标}}{\ph{说明}}\hfill
  \metric{\ph{指标}}{\ph{说明}}\hfill
  \metric{\ph{指标}}{\ph{说明}}
\end{frame}

% ============================ 3. 相关工作 ============================
\renewcommand{\secblurb}{\ph{本节一句话主旨：相对已有工作的关键差异}}
\section{相关工作}

\begin{frame}{研究领域与相关工作}
  \begin{columns}[T]
    \begin{column}{\zonehalf}
      \begin{block}{所属领域}
        \begin{itemize}
          \item \keyword{\ph{方向一}}
          \item \keyword{\ph{方向二}}
          \item \ph{方向三}
        \end{itemize}
      \end{block}
    \end{column}
    \begin{column}{\zonehalf}
      \begin{block}{与既有工作的关系}
        \begin{itemize}
          \item 基础：\shl{\ph{基线 / 算法}}
          \item 对比：\shl{\ph{对比方法}}
          \item 最相关：\shl{\ph{最相关工作}} —— \ph{差异}
        \end{itemize}
      \end{block}
    \end{column}
  \end{columns}
  \tightgap
  \begin{callout}[关键差异]
    \ph{本文相对最相关工作的一句话差异}。
  \end{callout}
\end{frame}

% ============================ 4. 方法 ============================
\renewcommand{\secblurb}{\ph{本节一句话主旨：方法的核心思想与组成}}
\section{方法}

\begin{frame}{方法 —— 总览}
  \begin{columns}[T]
    \begin{column}{\zonehalf}
      \begin{itemize}
        \item \shl{\ph{组件一}}：\ph{作用}。
        \item \shl{\ph{组件二}}：\ph{作用}。
        \item \ph{数据流 / 训练流程要点}。
      \end{itemize}
    \end{column}
    \begin{column}{\zonehalf}
      \centering
      \phfig
      \tightgap
      \figcap{2}{\ph{系统 / 方法框图}}
    \end{column}
  \end{columns}
\end{frame}

\begin{frame}{方法 —— 核心设计要素}
  \begin{itemize}
    \item \shl{\ph{要素 1}}：\ph{说明} $\to$ \ph{收益}。
    \item \shl{\ph{要素 2}}：\ph{说明}。
    \item \shl{\ph{要素 3}}：\ph{说明}。
    \item \shl{\ph{要素 4}}：\ph{说明}。
  \end{itemize}
  \tightgap
  \begin{callout}[要点]
    \ph{方法层面最重要的一句话结论}。
  \end{callout}
\end{frame}

% ============================ 5. 实验设计 ============================
\renewcommand{\secblurb}{\ph{本节一句话主旨：任务范围与评测协议}}
\section{实验设计}

\begin{frame}{实验设计}
  \begin{columns}[T]
    \begin{column}{\zonetext}
      \begin{block}{任务 / 数据}
        \begin{itemize}
          \setlength{\itemsep}{0.3em}
          \item \ph{任务一}
          \item \ph{任务二}
          \item \ph{任务三}
        \end{itemize}
      \end{block}
      \begin{exampleblock}{评测协议}
        \small \ph{评测方式、指标、试验次数}。\\[0.15em]
        {\footnotesize\color{sustechgrey} 基线：\ph{baseline A · B · C}}
      \end{exampleblock}
    \end{column}
    \begin{column}{\zonefig}
      \centering
      \phfig
      \tightgap
      \figcap{3}{\ph{任务图示}}
    \end{column}
  \end{columns}
\end{frame}

% ============================ 6. 实验结果 ============================
\renewcommand{\secblurb}{\ph{本节一句话主旨：最重要的量化结论}}
\section{实验结果}

\begin{frame}{实验结果 —— 主表}
  \centering
  \renewcommand{\arraystretch}{1.15}
  \begin{tabular}{lcccc}
    \toprule
    \theadrow \thc{\ph{任务}} & \thc{\ph{基线}} & \thc{\ph{本文}} & \thc{\ph{指标A}} & \thc{\ph{指标B}}\\
    \midrule
    \ph{行 1}        & \ph{·} & \shl{\ph{·}} & \ph{·} & \shl{\ph{·}}\\
    \altrow \ph{行 2} & \ph{·} & \shl{\ph{·}} & \ph{·} & \shl{\ph{·}}\\
    \ph{行 3}        & \ph{·} & \shl{\ph{·}} & \ph{·} & \shl{\ph{·}}\\
    \midrule
    \theadrow \thc{平均} & \thc{\ph{·}} & \thc{\ph{·}} & \thc{\ph{·}} & \thc{\ph{·}}\\
    \bottomrule
  \end{tabular}
  \deckgap
  \begin{callout}[结论]
    \ph{用一句话点出主表的核心结论（含关键数字）}。
  \end{callout}
\end{frame}

\begin{frame}{实验结果 —— 分析}
  \begin{columns}[T]
    \begin{column}{\zonehalf}
      \centering
      \phfig
      \tightgap
      \figcap{4}{\ph{曲线 / 分析图}}
    \end{column}
    \begin{column}{\zonehalf}
      \begin{itemize}
        \item \shl{\ph{现象一}}：\ph{解读}。
        \item \shl{\ph{现象二}}：\ph{解读}。
        \item \shl{\ph{现象三}}：\ph{解读}。
      \end{itemize}
    \end{column}
  \end{columns}
\end{frame}

% ============================ 7. 局限性 ============================
\renewcommand{\secblurb}{\ph{本节一句话主旨：主要局限与未来方向}}
\section{局限性}

\begin{frame}{局限性}
  \begin{itemize}
    \item \shl{\ph{局限一}}：\ph{说明}。
    \item \shl{\ph{局限二}}：\ph{说明}。
    \item \shl{\ph{局限三}}：\ph{说明}。
  \end{itemize}
  \begin{callout}[作者展望]
    \ph{作者给出的未来方向}。
  \end{callout}
\end{frame}

% ============================ 8. 要点总结 ============================
\renewcommand{\secblurb}{\ph{本节一句话主旨：全文最值得带走的结论}}
\section{要点总结}

\begin{frame}{要点总结（Takeaways）}
  \begin{enumerate}
    \item \ph{要点一}。
    \item \shl{\ph{要点二}}。
    \item \ph{要点三}。
  \end{enumerate}
  \begin{callout}[一句话带走]
    \ph{一句话总结全文}。
  \end{callout}
\end{frame}

% ============================ 9. 参考文献 ============================
\renewcommand{\secblurb}{\ph{本节一句话主旨：关键参考与根基工作}}
\section{参考文献}

\begin{frame}{参考文献}
  \setbeamertemplate{bibliography item}[text]
  \footnotesize
  \begin{thebibliography}{9}
    \bibitem{ref1} \ph{作者}. \textit{\ph{标题}}. \ph{会议/期刊}, \ph{年份}.（本文）
    \bibitem{ref2} \ph{作者}. \textit{\ph{标题}}. \ph{会议/期刊}, \ph{年份}.
    \bibitem{ref3} \ph{作者}. \textit{\ph{标题}}. \ph{会议/期刊}, \ph{年份}.
    \bibitem{ref4} \ph{作者}. \textit{\ph{标题}}. \ph{会议/期刊}, \ph{年份}.
  \end{thebibliography}
\end{frame}

% ============================ Q&A ============================
\begin{frame}[plain,c]
  \centering
  {\fontsize{40}{44}\selectfont\bfseries\color{emph} Q\,\&\,A}

  \deckgap
  {\Large \ph{致谢语}}

  \deckgap
  {\color{sustechgrey}\footnotesize \ph{汇报人} \textperiodcentered\ \ph{机构}\\
   \ph{联系方式 / 主页}}
\end{frame}

% ============================ 附录 / Backup ============================
\appendix

\begin{frame}[plain,c]
  \centering
  {\color{sustechgrey}\large Appendix}\\[0.6em]
  {\Huge\bfseries 备用幻灯片 / Backup}\\[0.9em]
  {\color{emph}\rule{2cm}{2pt}}
\end{frame}

\begin{frame}{附录 A —— \ph{补充主题}}
  \begin{columns}[T]
    \begin{column}{\zonetext}
      \begin{itemize}
        \item \shl{\ph{补充点一}}：\ph{说明}。
        \item \ph{补充点二}。
      \end{itemize}
    \end{column}
    \begin{column}{\zonefig}
      \centering
      \phfig
      \tightgap
      \figcap{5}{\ph{补充图}}
    \end{column}
  \end{columns}
  \tightgap
  \begin{callout}[备注]
    \ph{补充说明}。
  \end{callout}
\end{frame}

\end{document}
